\section{Plataforma Docker}

\begin{frame}
  \frametitle{O que é o Docker?}
  \begin{block}{Definição}
    Plataforma aberta para \textbf{desenvolvimento, envio e execução} de aplicações.
  \end{block}

  \vspace{0.5em}
  \begin{itemize}
    \item Separa a aplicação de sua infraestrutura
    \item Permite desenvolver software de forma mais rápida
  \end{itemize}
\end{frame}


\begin{frame}
  \frametitle{Containers}
  \begin{block}{O que é um container?}
    Ambiente \textbf{isolado} que empacota tudo o que a aplicação precisa para rodar sem depender do que está instalado no \textit{host}.
  \end{block}

  \vspace{0.5em}
  \begin{itemize}
    \item Vários containers rodando simultaneamente com segurança
    \item Leves e portáveis
    \item Fáceis de compartilhar em equipe
  \end{itemize}
\end{frame}

\begin{frame}
  \frametitle{Fluxo de Trabalho com Containers}
  O Docker provê ferramentas e uma plataforma para gerenciar os containers, seguindo o fluxo:

  \vspace{0.5em}
  \begin{enumerate}
    \item<2-> \textbf{Desenvolver} a aplicação e seus componentes usando containers
    \item<3-> O container vira a unidade de \textbf{distribuição e teste}
    \item<4-> \textbf{Implantar} em produção como container ou serviço
  \end{enumerate}
\end{frame}

\begin{frame}
  \frametitle{Ciclo de Desenvolvimento}
  \begin{columns}
    \column{0.5\textwidth}
    \textbf{Desenvolvimento}
    \begin{itemize}
      \item Código local em containers
      \item Compartilhamento com a equipe
    \end{itemize}

    \column{0.5\textwidth}
    \textbf{Entrega}
    \begin{itemize}
      \item Testes manuais e/ou automáticos
      \item Correção de bugs no próprio ambiente
      \item \textit{Push} da imagem atualizada para produção
    \end{itemize}
  \end{columns}
\end{frame}

\begin{frame}
  \frametitle{Eficiência de Recursos}
  \begin{alertblock}{Docker vs Máquinas Virtuais}
    Docker é leve e rápido, com melhor custo-benefício que VMs baseadas em \textit{hypervisor}.
  \end{alertblock}
\end{frame}